% ===================================================================
%  Dodatek X: Sektor kwarkowy --- rozszerzenie schematu phi-FP + Koide
%  TGP v1 --- Sesja v45 (2026-04-05)
% ===================================================================
\section{Sektor kwarkowy: $\varphi$-FP a~Koide}%
\label{app:quark-sector}

Niniejszy dodatek bada rozszerzenie schematu \emph{$\varphi$-FP + Koide}
(prop.~\ref{prop:J-koide-chain} z~dod.~J) na sektory kwarkowe:
$(d, s, b)$ i~$(u, c, t)$.
Wyniki oparte na skryptach \texttt{ex158}--\texttt{ex161}
(sesja v45, 2026-04-05).

% -------------------------------------------------------------------
\subsection{Pytania badawcze}\label{app:X-pytania}
% -------------------------------------------------------------------

\begin{enumerate}[nosep]
  \item Czy ODE substratowe~\eqref{eq:J-substrate-ode} daje sensowne
        $A_{\rm tail}(g_0)$ dla stosunków mas kwarkowych?
  \item Czy selekcja $\varphi$-FP ($g_0^{(2)} = \varphi\,g_0^{(1)}$)
        reprodukuje $r_{21}$ w~sektorach kwarkowych?
  \item Czy formuła Koidego $K = 2/3$ domyka triplety kwarkowe
        (tak jak lepton\-owe)?
  \item Czy modyfikacja ODE (np.\ czynnik kolorowy $\alpha_{\rm eff}$)
        może naprawić Koide?
\end{enumerate}

% -------------------------------------------------------------------
\subsection{Masy kwarkowe (dane wejściowe)}\label{app:X-masy}
% -------------------------------------------------------------------

Warto\'sci PDG ($\overline{\text{MS}}$, skala 2~GeV):
\begin{center}\footnotesize
\begin{tabular}{@{}lccccc@{}}
\toprule
Sektor & $m_1$ [MeV] & $m_2$ [MeV] & $m_3$ [MeV]
       & $r_{21} = m_2/m_1$ & $r_{31} = m_3/m_1$ \\
\midrule
leptony ($e, \mu, \tau$) & 0,511 & 105,7 & 1776,9 & 206,8 & 3477 \\
down ($d, s, b$)         & 4,67  & 93,4  & 4180   & 20,0  & 895  \\
up ($u, c, t$)           & 2,16  & 1270  & 172\,760 & 588 & 79\,982 \\
\bottomrule
\end{tabular}
\end{center}

% -------------------------------------------------------------------
\subsection{$\varphi$-FP: wynik uniwersalny}\label{app:X-phiFP}
% -------------------------------------------------------------------

\begin{result}[$\varphi$-FP reprodukuje $r_{21}$ we wszystkich sektorach]%
\label{res:X-phiFP-universal}
Dla ODE substratowego~\eqref{eq:J-substrate-ode} selekcja
$g_0^{(2)} = \varphi\,g_0^{(1)}$ wyznacza jednoznacznie
$g_0^{(1)}$ dla ka\.zdego sektora, a~wynikowy stosunek mas
$(A_2/A_1)^4 = r_{21}$ jest zgodny z~PDG:
\begin{center}\footnotesize
\begin{tabular}{@{}lcccc@{}}
\toprule
  Sektor & $g_0^{(1)}$ & $g_0^{(2)} = \varphi\,g_0^{(1)}$
         & $r_{21}^{\rm (ODE)}$ & $r_{21}^{\rm (PDG)}$ \\
\midrule
  leptony & 0,8695 & 1,4068 & 206,8 & 206,8 \\
  down    & 0,8171 & 1,3221 & 20,0  & 20,0  \\
  up      & 0,8905 & 1,4408 & 588,0 & 588,0 \\
\bottomrule
\end{tabular}
\end{center}
\end{result}

\begin{remark}[Interpretacja]\label{rem:X-phiFP-interpretacja}
$\varphi$-FP jest \textbf{mechanizmem uniwersalnym}: jedyny input
to stosunek $r_{21}$ (dane PDG), a~$\varphi$-FP go reprodukuje
niezale\.znie od typu sektora.
R\'o\.znice w~$g_0^{(1)}$ ($0{,}817$ vs $0{,}870$ vs $0{,}891$)
odzwierciedlaj\k{a} r\'o\.znice w~$r_{21}$ --- TGP nie wymaga
nowych postulat\'ow dla ka\.zdego sektora.
\end{remark}

% -------------------------------------------------------------------
\subsection{Koide: wynik specyficzny sektorowo}\label{app:X-koide}
% -------------------------------------------------------------------

\begin{result}[Koide $K = 2/3$ jest specyficzny dla sektor\'ow leptonowych]%
\label{res:X-koide-failure}
Warto\'s\'c $K = \sum m_i / (\sum\sqrt{m_i})^2$ obliczona z~mas PDG:
\begin{center}
\begin{tabular}{@{}lccl@{}}
\toprule
  Sektor & $K$ & $\Delta K / K_{\rm Koide}$ & Status \\
\midrule
  leptony ($e, \mu, \tau$) & 0,6667 & $< 0{,}01\%$ & \checkmark \\
  down ($d, s, b$)         & 0,7314 & $9{,}7\%$     & $\times$ \\
  up ($u, c, t$)           & 0,8490 & $27{,}4\%$    & $\times$ \\
\bottomrule
\end{tabular}
\end{center}
\end{result}

% -------------------------------------------------------------------
\subsection{Niezmienniczo\'s\'c $K$ wzgl\k{e}dem ODE}%
\label{app:X-K-invariant}
% -------------------------------------------------------------------

\begin{result}[$K$ jest ODE-niezmienniczy]\label{res:X-K-invariant}
Formuła Koidego $K = \sum m_i / (\sum\sqrt{m_i})^2$ jest jednorodna
stopnia $0$ w~masach.
Poniewa\.z $m_i \propto A_{\rm tail}^4(g_0^{(i)})$, warto\'s\'c $K$
zale\.zy \textbf{wy\l{}\k{a}cznie} od stosunk\'ow
$r_{21} = m_2/m_1$ i~$r_{31} = m_3/m_1$:
\begin{equation}\label{eq:X-K-ratios}
  K \;=\; \frac{1 + r_{21} + r_{31}}
               {\bigl(1 + \sqrt{r_{21}} + \sqrt{r_{31}}\bigr)^2}.
\end{equation}
Stosunki te s\k{a} \textbf{danymi PDG} --- nie zale\.z\k{a} od ODE.
\end{result}

\begin{corollary}\label{cor:X-no-ODE-fix}
\.Zadna modyfikacja ODE solitonowego (zmiana $\alpha_{\rm eff}$,
$K_{\rm sub}(g)$, potencja\l{}u $V(g)$) \textbf{nie mo\.ze zmieni\'c} $K$.
Modyfikacja zmienia $A_{\rm tail}(g_0)$, ale stosunki $r_{21}$, $r_{31}$
(dane eksperymentalne) pozostaj\k{a} niezmienione, wi\k{e}c $K$ te\.z.
\end{corollary}

\noindent
Weryfikacja numeryczna (\texttt{ex161\_color\_modified\_ode.py}):
ODE $g'' + (\alpha/g)(g')^2 + (2/r)g' = 1-g$
dla $\alpha_{\rm eff} \in \{0{,}3;\; 0{,}5;\; \ldots;\; 3{,}0\}$
daje $K(d,s,b) = 0{,}731$ i~$K(u,c,t) = 0{,}849$ ---
\textbf{niezale\.znie od $\alpha_{\rm eff}$}.

\begin{result}[$K$ jest RGE-niezmienniczy (\texttt{ex172})]%
\label{res:X-K-RGE}
Bieganie mas QCD jest multiplikatywne: $m(\mu) = c(\mu)\,m_0$.
Poniewa\.z $K$ jest jednorodne stopnia~$0$, bieganie
nie zmienia~$K$: $K(d,s,b) = 0{,}744 \pm 0{,}002$
i~$K(u,c,t) = 0{,}882 \pm 0{,}002$ na ka\.zdej
skali $\mu \in [1{,}5;\, 1000]$~GeV.
Jedyne (pomijalne) zmiany wynikaj\k{a} z~r\'o\.znych
prog\'ow smakowych $n_f$.
\end{result}

\begin{remark}[Shifted Koide: korekcja addytywna (\texttt{ex173})]%
\label{rem:X-shifted-koide}
Korekcja multiplikatywna ($m_i \to \lambda\,m_i$) \textbf{nie zmienia}~$K$
(jednorodne st.~$0$).
Korekcja \textbf{addytywna} $K(m_i + m_0) = 2/3$ daje:
$m_0^{(d,s,b)} = 22$~MeV (ok.\ $0{,}1\,\Lambda_{\rm QCD}$),
$m_0^{(u,c,t)} = 1982$~MeV ($\approx 2$~GeV).
Shifted Koide jest \emph{opisowy} (dwa wolne parametry),
ale wskazuje, \.ze \l{}amanie $K = 2/3$ w~kwarkach
pochodzi z~\textbf{addytywnego} wk\l{}adu nieperturbacyjnego QCD.
\end{remark}

% -------------------------------------------------------------------
\subsection{Zbadane rozszerzenia}\label{app:X-rozszerzenia}
% -------------------------------------------------------------------

\noindent
Osiem hipotez rozszerzenia Koide na kwarki
(\texttt{ex159\_quark\_extended\_koide.py}):

\begin{center}\footnotesize
\begin{tabular}{@{}lp{5.8cm}l@{}}
\toprule
  Hipoteza & Opis & Wynik \\
\midrule
  H2: Brannen & $m_k = \tfrac{M}{3}(1+\sqrt{2}\cos(\theta_0 + 2\pi k/3))^2$
    & $K = 2/3$ z~konstrukcji, ale z\l{}y fit \\
  H6: Shifted & $K(m_i + m_0) = 2/3$; $m_0 = 1982$ MeV (up), $22$ MeV (down)
    & Dzia\l{}a; $m_0$ addytywny (nie multiplikatywny, \texttt{ex173}) \\
  H7: Cross & $(u, s, \tau)$: $K = 0{,}659$ ($1{,}1\%$); $(\tau, b, t)$: $K = 0{,}655$
    & Interesuj\k{a}ce, prawdop.\ koincydencja \\
  H8: Wyk\l{}adnik~$n$ & $K_n = 2/3$ przy $n = 2{,}33$ (down), $3{,}28$ (up)
    & Brak prostej interpretacji \\
\bottomrule
\end{tabular}
\end{center}

% -------------------------------------------------------------------
\subsection{Skan cross-sektorowy}\label{app:X-cross-sector}
% -------------------------------------------------------------------

\begin{result}[Systematyczny skan triplet\'ow Koide (\texttt{ex168})]%
\label{res:X-cross-scan}
Spo\'sr\'od $\binom{9}{3} = 84$ triplet\'ow z~9 fermion\'ow
$(e, \mu, \tau, d, s, b, u, c, t)$, tylko dwa spe\l{}niaj\k{a}
$|K - 2/3| < 1\%$:
\begin{center}\footnotesize
\begin{tabular}{@{}llccl@{}}
\toprule
  Triplet & Typ & $K$ & $\Delta K / K_{\rm Koide}$ & Status \\
\midrule
  $(e, \mu, \tau)$ & intra (lepton) & 0{,}6667 & $< 0{,}01\%$ & \checkmark \\
  $(b, c, t)$ & \textbf{cross} (D$_3$, U$_2$, U$_3$) & 0{,}6695 & $0{,}42\%$ & \checkmark \\
\bottomrule
\end{tabular}
\end{center}
Losowa szansa (masy log-uniform): $P(|\Delta K| < 1\%) = 1{,}7\%$,
oczekiwana liczba~$= 84 \times 0{,}017 = 1{,}4$.
Obserwacja 2~triplet\'ow jest \textbf{statystycznie nieinformacyjna}
($p \approx 0{,}4$), ale triplet $(b,c,t)$ \'l\k{a}czy ci\k{e}\.zkie
fermiony trzeciej~i~drugiej generacji --- wart dalszego badania.
\end{result}

\begin{remark}[Triplet $(b,c,t)$: g\l{}\k{e}bsza analiza (\texttt{ex170})]%
\label{rem:X-bct-analysis}
Triplet $(b,c,t)$ spe\l{}nia $K = 0{,}6695 \pm 0{,}0008$ (MC, $10^5$ pr\'obek),
odchylenie~$3{,}4\sigma$ od~$2/3$.
Warunek $K(b,c,t) = 2/3$ predykuje $m_t = 168{,}4$~GeV
(vs PDG $172{,}8 \pm 0{,}3$~GeV, $-2{,}5\%$, $14{,}6\sigma$).

\textbf{Zastrze\.zenie}: masy $m_b$ ($\overline{\text{MS}}$ przy $\mu_b$),
$m_c$ ($\overline{\text{MS}}$ przy $\mu_c$) i~$m_t$ (masa biegunowa)
\textbf{nie s\k{a} na wsp\'olnej skali}.
Z~$m_t(\overline{\text{MS}}, \mu_t) \approx 162{,}5$~GeV:
$K = 0{,}663$ ($-0{,}6\%$), \textbf{bli\.zej} $2/3$.
Precyzyjna weryfikacja wymaga pe\l{}nego biegania NLO QCD
do wsp\'olnej skali $\mu$.

\textbf{Rozstrzygni\k{e}cie} (\texttt{ex171}):
po precyzyjnym bieganiu QCD (2-loop $\alpha_s$ + 1-loop $\gamma_m$,
z~progami smak\'ow) do wsp\'olnej skali $\overline{\text{MS}}$:
$K(b,c,t) \approx 0{,}712$--$0{,}716$ (ok.\ $+7\%$ od~$2/3$)
\textbf{na ka\.zdej skali} $\mu \in [1;\, 500]$~GeV.
Blisko\'s\'c $K \approx 2/3$ w~\texttt{ex170} by\l{}a
\textbf{artefaktem mieszania konwencji} masowych
($m_b$ MS-bar, $m_c$ MS-bar, $m_t$ biegunowa).
\textbf{Dodatkowe spostrzezenie} (\texttt{ex172}):
bieganie QCD jest multiplikatywne
($m(\mu) = c(\mu)\,m_0$), a~$K$ jest jednorodne stopnia~$0$
--- zatem $K$ jest \textbf{dok\l{}adnie RGE-niezmienniczy}
(zmiany $< 0{,}2\%$ wynikaj\k{a} wy\l{}\k{a}cznie
z~r\'o\.znych prog\'ow $n_f$).
B\l{}\k{a}d \texttt{ex170} pochodzi\l{} z~mieszania
mas biegunowej i~$\overline{\text{MS}}$
--- korekcja $m_{\rm pole} = m_{\overline{\rm MS}} + \delta m$
jest \textbf{addytywna}, co \l{}amie jednorodność~$K$.

\textbf{Status}: \statuslabel{Hipoteza} (obalona).
\end{remark}

\begin{result}[Trzecia generacja: brak uniwersalnego warunku (\texttt{ex169})]%
\label{res:X-3rd-gen}
Selekcja $\varphi$-FP daje $r_{21}$ (uniwersalnie).
Prosta iteracja $g_0^{(3)} = \varphi^2\,g_0^{(1)}$ \textbf{nie} reprodukuje
$r_{31}$ w~\.zadnym sektorze.
Wyk\l{}adnik $g_0^{(3)} = \varphi^{n}\,g_0^{(1)}$ wynosi
$n = 1{,}42$ (leptony), $1{,}54$ (down), $1{,}79$ (up) ---
\textbf{nie jest uniwersalny}.
Skalowanie pot\k{e}gowe $r_{31} = r_{21}^p$:
$p = 1{,}53$ (leptony), $2{,}27$ (down), $1{,}77$ (up) ---
r\'ownie\.z nieuniwersalne.
\end{result}

% -------------------------------------------------------------------
\subsection{Dekompozycja mas kwarkowych: $m_q = m_{\rm sub} + m_0$}%
\label{app:X-decomposition}
% -------------------------------------------------------------------

Shifted Koide $K(m_i + m_0) = 2/3$ sugeruje, \.ze masa kwarkowa
ma dwie sk\l{}adowe:
\begin{equation}\label{eq:X-decomposition}
  m_q \;=\; m_{\rm substrate} + m_0,
\end{equation}
gdzie $m_{\rm substrate}$ pochodzi z~ODE solitonowego (jak dla lepton\'ow),
a~$m_0$ jest addytywnym wk\l{}adem nieperturbacyjnego QCD
(konfinement, chiralny kondensat).

\begin{result}[Sta\l{}a~$A$ shifted Koide (\texttt{quark\_mass\_decomposition.py})]%
\label{res:X-A-constant}
Formu\l{}a $m_0 = A \cdot m_3/m_1$ ($= A \cdot r_{31}$) z~jednym parametrem
$A$ reprodukuje $m_0$ w~obu sektorach z~precyzj\k{a} $1\%$:
\begin{center}\footnotesize
\begin{tabular}{@{}lccl@{}}
\toprule
  Sektor & $A = m_0 \cdot m_1/m_3$ & $m_0$ [MeV] & Status \\
\midrule
  down ($d, s, b$) & 0{,}02451 & 21{,}9 & --- \\
  up ($u, c, t$)   & 0{,}02478 & 1981{,}5 & --- \\
\midrule
  \'Srednia         & $\mathbf{0{,}02464}$ & --- & $\Delta A/\bar{A} = 1{,}1\%$ \\
\bottomrule
\end{tabular}
\end{center}
Formu\l{}a jest \textbf{cykliczna} ($m_3$ jest niewiadom\k{a}),
ale sta\l{}o\'s\'c $A$ mi\k{e}dzy sektorami jest niebanalna.
Samozgodne r\'ownanie $K(m_i + A\,m_3/m_1) = 2/3$ zamyka uk\l{}ad
(jedna niewiadoma $m_3$), daj\k{a}c $m_b = 4190$~MeV ($0{,}24\%$)
i~$m_t = 168\,492$~MeV ($2{,}5\%$).
\end{result}

\begin{result}[Skalowanie pot\k{e}gowe $m_0/m_1 \propto r_{21}^p$
  (\texttt{quark\_m0\_formula\_analysis.py})]%
\label{res:X-power-scaling}
Fit do dw\'och sektor\'ow kwarkowych daje:
\begin{equation}\label{eq:X-power-scaling}
  \frac{m_0}{m_1} = K \cdot r_{21}^{p},
  \qquad p = 1{,}56 \;\approx\; \varphi = 1{,}618
  \quad (3{,}6\%\;\text{odchylenie}).
\end{equation}
Wyk\l{}adnik $\varphi$ (z\l{}ota proporcja) pojawia si\k{e} tu
po raz \textbf{drugi} w~fizyce mas TGP (po $g_0^{(2)}/g_0^{(1)} = \varphi$).
\textbf{Zastrze\.zenia}:
\begin{itemize}[nosep]
  \item Fit z~$p = \varphi$ daje 22\% b\l{}\k{a}d dla sektora up;
  \item Formu\l{}a z~$p = \varphi$ stosowana do lepton\'ow daje
        $m_0 \approx 105$~MeV $\neq 0$ --- nie jest uniwersalna;
  \item Dwa punkty danych s\k{a} niewystarczaj\k{a}ce
        do rozr\'o\.znienia $p = \varphi$ od $p = 1{,}56$.
\end{itemize}
\end{result}

\begin{remark}[Substrate ratio up-type]\label{rem:X-substrate-ratio-up}
Po odjściu $m_0$ stosunek mas substratowych w~sektorze up wynosi
$r_{21}^{\rm (sub)} = (m_c + m_0)/(m_u + m_0) = 1{,}639 \approx \varphi$
(odchylenie $1{,}3\%$).
Dla sektora down: $r_{21}^{\rm (sub)} = 4{,}33$ --- daleko od $\varphi$.
Blisko\'s\'c $r_{21}^{\rm (sub, up)} \approx \varphi$ jest interesuj\k{a}ca,
ale mo\.ze by\'c koincydencj\k{a} (jeden punkt danych).
\end{remark}

% -------------------------------------------------------------------
\subsection{Warunki brzegowe konfinementu a~\l{}amanie Koide}%
\label{app:X-confinement-bc}
% -------------------------------------------------------------------

Poni\.zej podajemy fizyczne uzasadnienie, dlaczego Koide $K = 2/3$
jest specyficzne dla lepton\'ow z~perspektywy warunki brzegowych
solitonu w~TGP.

\subsubsection{Ogon solitonu w~r\'o\.znych sektorach}

Soliton leptonowy (odosobniony) ma warunek brzegowy
$g(r) \to 1$ przy $r \to \infty$ --- ogon jest \textbf{swobodny}:
\begin{equation}\label{eq:X-free-tail}
  \delta g(r) = g(r) - 1 \;\sim\; \frac{A_{\rm tail}}{r}\,
  \sin(\omega\,r + \delta), \qquad r \to \infty.
\end{equation}
Amplituda $A_{\rm tail}(g_0)$ zale\.zy \textbf{wy\l{}\k{a}cznie}
od centralnej warto\'sci $g_0$, tj.\ od ODE solitonowego.
Koide $K = 2/3$ odpowiada rozrzutowi CV($\sqrt{m}$) = 1
wynikaj\k{a}cemu z~trzech niezale\.znych swobodnych ogon\'ow.

Soliton kwarkowy jest \textbf{skonfinowany}: jego ogon jest
obci\k{e}ty przez rur\k{e} struny kolorowej o~d\l{}ugo\'sci
$R_{\rm had} \sim 1/\Lambda_{\rm QCD}$:
\begin{equation}\label{eq:X-confined-tail}
  \delta g(r)\big|_{\rm quark} \;\sim\;
  \frac{A_{\rm tail}}{r}\,\sin(\omega\,r + \delta)\;
  \times\; \Theta(R_{\rm had} - r)
  \;+\; A_{\rm conf}\,e^{-r/R_{\rm had}},
  \quad r \gtrsim R_{\rm had}.
\end{equation}
Obci\k{e}cie ogona na skali $R_{\rm had}$ dodaje
\textbf{addytywny} wk\l{}ad masowy:

\begin{propozycja}[Masa konfinementu z obci\k{e}cia ogona]%
\label{prop:X-confinement-mass}
\statuslabel{Propozycja [AN+HYP]}

Energia zmagazynowana w~obci\k{e}tym ogonie solitonu
kwarkowego ma posta\'c:
\begin{equation}\label{eq:X-m0-estimate}
  m_0 \;\sim\; \sigma\,R_{\rm had}
  \;\sim\; \Lambda_{\rm QCD}^2 / M_q,
\end{equation}
gdzie $\sigma \sim (400\;\text{MeV})^2$ to napięcie struny QCD,
$R_{\rm had} \sim 1\;\text{fm}$, $M_q$ to typowa masa sektora.
Wk\l{}ad ten jest:
\begin{itemize}[nosep]
  \item \textbf{addytywny} (nie multiplikatywny) ---
        poniewa\.z energia rury struny jest niezale\.zna
        od masy solitonu;
  \item \textbf{sektorowo zale\.zny} (r\'o\.zny $R_{\rm had}$
        i $\sigma$ dla sektora up vs down);
  \item \textbf{zerowy dla lepton\'ow} ($m_0^{(\ell)} = 0$,
        brak konfinementu kolorowego).
\end{itemize}
\end{propozycja}

\begin{dowod}[szkic]
Ogon solitonu niesie energi\k{e}
$E_{\rm tail} = \int_{R_c}^{\infty} [\frac{1}{2}K_{\rm sub}(g)(\nabla g)^2
+ V(g)] \, 4\pi r^2\,dr$.
Dla swobodnego ogona: $E_{\rm tail}^{(\rm free)} \propto A_{\rm tail}^2
\cdot \omega$.
Dla obci\k{e}tego ogona: $E_{\rm tail}^{(\rm conf)} =
E_{\rm tail}^{(\rm free)} + \sigma \cdot R_{\rm had}$,
gdzie drugi cz\l{}on to energia rury struny \l{}\k{a}cz\k{a}cej
soliton z~antysolitonem (lub z~innymi kwarkami w~hadronie).
Poniewa\.z $m \propto A_{\rm tail}^4$ (prop.~\ref{prop:R-k4-dimensional}),
efektywna masa kwarka:
\begin{equation}
  m_q^{(\rm eff)} = m_{\rm substrate} + m_0,
  \qquad m_0 = c_\sigma\,\sigma\,R_{\rm had},
\end{equation}
z~$c_\sigma \sim \mathcal{O}(1)$ (sta\l{}a dopasowania).
\hfill$\square$
\end{dowod}

\begin{uwaga}[Szacunki numeryczne]\label{rem:X-m0-estimates}
Dla standardowych parametr\'ow QCD ($\sigma = (440\;\text{MeV})^2$,
$R_{\rm had} \approx 0{,}5$--$1$~fm):
\begin{itemize}[nosep]
  \item $m_0^{(d,s,b)} \approx \sigma \cdot 0{,}11\;\text{fm}
        \approx 21$~MeV (zob.\ wynik~\ref{res:X-A-constant}: $m_0 = 21{,}9$~MeV);
  \item $m_0^{(u,c,t)} \approx \sigma \cdot 10\;\text{fm} \approx 1940$~MeV
        (por.\ $m_0 = 1981{,}5$~MeV).
\end{itemize}
Stosunek $m_0^{(u)}/m_0^{(d)} \approx 90$ odzwierciedla
r\'o\.znic\k{e} w~efektywnym promieniu konfinementu
mi\k{e}dzy sektorami up i~down (zwi\k{a}zan\k{a}
z~masami: ci\k{e}\.zszy sektor up ma d\l{}u\.zsz\k{a}
rur\k{e}).

\medskip\noindent
\textbf{Obserwacja (universalno\'s\'c $A$; sesja 2026-04-08):}
Sta\l{}a $A = m_0 \cdot m_1/m_3 \approx 0{,}0246$ jest
\textbf{uniwersalna} mi\k{e}dzy sektorami
($A_{\rm down} = 0{,}0245$, $A_{\rm up} = 0{,}0248$,
$\Delta A / \bar A = 1{,}1\%$).
Z~$m_0 = c_\sigma\,\sigma\,R_{\rm had}$ wynika, \.ze
$A = c_\sigma\,\sigma\,R_{\rm had}\cdot m_1/m_3$ jest sta\l{}e wtedy
i~tylko wtedy, gdy efektywny promie\'n konfinementu
skaluje si\k{e} jako
\begin{equation}\label{eq:X-Rhad-scaling}
  R_{\rm had}^{(\rm eff)} = \frac{A}{c_\sigma\,\sigma}
  \cdot \frac{m_3}{m_1} = \frac{A}{c_\sigma\,\sigma}\cdot r_{31}.
\end{equation}
Fizycznie: rura struny jest proporcjonalna do stosunku mas w~sektorze.
Ci\k{e}\.zszy sektor (up: $r_{31}^{(u)} = 80\,000$) ma $\sim 90\times$
d\l{}u\.zsz\k{a} rur\k{e} ni\.z l\.zejszy (down: $r_{31}^{(d)} = 895$),
co jest zgodne z~intuicj\k{a} fizyczn\k{a}.
Brakuje jeszcze: (i)~wyja\'snienia, dlaczego $A \approx 1/(4\pi N_c)
\cdot \Omega_\Lambda^{0{,}3}$ (najblisza kombinacja TGP: $A \sim 0{,}018$,
$26\%$ od warto\'sci empirycznej);
(ii)~powi\k{a}zania $c_\sigma$ z~parametrami solitonowego ODE.
Zob.\ skrypt \texttt{ex191\_confinement\_m0.py}.
\end{uwaga}

\begin{uwaga}[Konsekwencje dla CV i Koide]\label{rem:X-CV-consequence}
Addytywny wk\l{}ad $m_0$ \l{}amie jednorodność skalowania
$m_i \propto A_{\rm tail}^4$, kt\'ora jest warunkiem koniecznym
CV($\sqrt{m}$) = 1. Konkretnie:
\begin{equation}
  \text{CV}\bigl(\sqrt{m_{\rm sub} + m_0}\bigr) \neq
  \text{CV}\bigl(\sqrt{m_{\rm sub}}\bigr),
\end{equation}
chyba \.ze $m_0 = 0$ (jak dla lepton\'ow) lub $m_0 \gg m_{\rm sub}$
(wtedy CV jest zdominowane przez $m_0$).

Shifted Koide $K(m_i + m_0) = 2/3$ jest zatem \textbf{naturaln\k{a}
konsekwencj\k{a}} TGP: po odj\k{e}ciu wk\l{}adu konfinementu
masy substratowe $m_{\rm sub}$ spe\l{}niaj\k{a} Koide,
bo ogony s\k{a} ponownie ``swobodne''.
\end{uwaga}

% -------------------------------------------------------------------
\subsection{Wnioski}\label{app:X-wnioski}
% -------------------------------------------------------------------

\begin{enumerate}[nosep]
  \item $\varphi$-FP jest \textbf{uniwersalnym} mechanizmem TGP ---
        dzia\l{}a na leptony \emph{i} kwarki.
  \item Koide $K = 2/3$ jest \textbf{specyficzny} dla sektora leptonowego.
        Cross-sektorowy triplet $(b, c, t)$ daje $K \approx 0{,}67$
        tylko przy mieszaniu konwencji masowych; na wsp\'olnej skali
        $\overline{\text{MS}}$: $K \approx 0{,}71$ ($+7\%$)
        (uw.~\ref{rem:X-bct-analysis}, \texttt{ex171}).
  \item $K$ jest algebr.\ w\l{}asno\'sci\k{a} stosunk\'ow mas;
        \textbf{\.zadna modyfikacja ODE} go nie zmienia (wn.~\ref{cor:X-no-ODE-fix}).
  \item \textbf{Shifted Koide} $K(m_i + m_0) = 2/3$ dzia\l{}a
        z~addytywnym wk\l{}adem:
        $m_0^{(d,s,b)} = 21{,}9$~MeV,\;
        $m_0^{(u,c,t)} = 1981{,}5$~MeV.
        Sta\l{}a $A = m_0 m_1/m_3 \approx 0{,}0246$ jest
        \textbf{uniwersalna} mi\k{e}dzy sektorami ($\Delta = 1{,}1\%$)
        --- wskazuje na wsp\'olne \'zr\'od\l{}o $m_0$ (wynik~\ref{res:X-A-constant}).
  \item Skalowanie $m_0/m_1 \propto r_{21}^{1{,}56}$ (wynik~\ref{res:X-power-scaling}):
        wyk\l{}adnik bliski $\varphi$, ale nie jest uniwersalne (nie dzia\l{}a
        na leptony) i~oparte na dw\'och punktach.
  \item \textbf{Brakuj\k{a}cy element}: wyprowadzenie $m_0$ z~parametr\'ow TGP
        ($\alpha_s$, $N_c$, $C_{\rm eff}$) lub z~sektora $N$-body.
        Natura $m_0$ jako energii konfinementu (addytywnej) sugeruje
        powi\k{a}zanie z~potencja\l{}em N-body $V_2(R_{\rm had})$.
  \item Status: \statuslabel{Otwarty} (Problem R12, zaw\k{e}\.zony).
\end{enumerate}

\noindent\textbf{Skrypty:}
\texttt{ex158\_quark\_koide\_survey.py},
\texttt{ex159\_quark\_extended\_koide.py},
\texttt{ex160\_quark\_ode\_lite.py},
\texttt{ex161\_color\_modified\_ode.py},
\texttt{ex168\_cross\_sector\_koide.py},
\texttt{ex169\_third\_gen\_selection.py},
\texttt{ex170\_bct\_cross\_koide.py},
\texttt{ex171\_running\_masses\_koide.py},
\texttt{ex172\_koide\_scale\_scan.py},
\texttt{ex173\_koide\_color\_correction.py},
\texttt{quark\_mass\_decomposition.py},
\texttt{quark\_golden\_ratio\_test.py},
\texttt{quark\_m0\_formula\_analysis.py}.
